\documentclass[11pt]{article}
\usepackage[T1]{fontenc}
\usepackage[utf8]{inputenc}
\usepackage{lmodern}
\usepackage{microtype}
\usepackage[a4paper,margin=29mm]{geometry}
\usepackage{amsmath,amssymb,amsthm,mathtools}
\usepackage{booktabs,tabularx,longtable,array}
\usepackage{enumitem}
\usepackage{xcolor}
\usepackage{hyperref}
\usepackage[nameinlink,noabbrev]{cleveref}
\usepackage{fancyhdr}
\usepackage{lastpage}
\usepackage{verbatim}

\hypersetup{
  colorlinks=true,
  linkcolor=blue!45!black,
  citecolor=green!35!black,
  urlcolor=blue!55!black,
  pdftitle={Reciprocal-Power Sign Rigidity and Stieltjes Moment Lattices},
  pdfauthor={Research progress report}
}
\setlist{nosep,leftmargin=2.1em}
\setlength{\parindent}{0pt}
\setlength{\parskip}{0.55em}
\allowdisplaybreaks
\pagestyle{fancy}
\fancyhf{}
\lhead{Alaoglu--Erd\H{o}s progress report}
\rhead{\thepage/\pageref{LastPage}}
\renewcommand{\headrulewidth}{0.3pt}

\newtheorem{theorem}{Theorem}[section]
\newtheorem{proposition}[theorem]{Proposition}
\newtheorem{lemma}[theorem]{Lemma}
\newtheorem{corollary}[theorem]{Corollary}
\newtheorem{criterion}[theorem]{Criterion}
\theoremstyle{definition}
\newtheorem{definition}[theorem]{Definition}
\newtheorem{problem}[theorem]{Open problem}
\newtheorem{experiment}[theorem]{Exact experiment}
\theoremstyle{remark}
\newtheorem{remark}[theorem]{Remark}

\newcommand{\N}{\mathbb N}
\newcommand{\Z}{\mathbb Z}
\newcommand{\Q}{\mathbb Q}
\newcommand{\R}{\mathbb R}
\newcommand{\dd}{\,\mathrm d}
\newcommand{\cst}{c_{\theta}}
\newcommand{\UU}{\mathcal U}
\newcommand{\VV}{\mathcal V}
\newcommand{\EE}{\mathcal E}
\newcommand{\Sm}{\mathcal S}
\newcommand{\Pow}{\mathcal P}
\newcommand{\DeltaU}{\Delta(\UU)}
\newcommand{\AE}{Alaoglu--Erd\H{o}s}
\DeclareMathOperator{\sgn}{sgn}
\DeclareMathOperator{\diag}{diag}
\DeclareMathOperator{\Res}{Res}
\DeclareMathOperator{\sRes}{sRes}
\DeclareMathOperator{\Span}{span}

\title{\textbf{Reciprocal-Power Sign Rigidity and Stieltjes Moment Lattices}\\[0.4em]
\large A nonredundant continuation toward the \AE\ conjecture}
\author{Research progress report}
\date{22 August 2026}

\begin{document}
\maketitle

\begin{abstract}
Let $x\in\R$ satisfy $2^x,3^x\in\Z$.  The conjecture asserts that $x\in\Z$.
This report continues the attached unified manuscript without reproducing its surveys of
standard determinant, $p$-adic, interpolation, automaticity, or point-counting approaches.
The new work starts from the reciprocal exponent $\theta=1/x$ and the integer semigroup
correspondence
\[
  A^mB^n \longmapsto 2^m3^n,
  \qquad A=2^x,\quad B=3^x.
\]
Four concrete advances are obtained.  First, every counterexample gives an infinite
hierarchy of strict signs of explicitly integral generalized Vandermonde determinants;
a wrong sign is therefore a finite, logarithm-free contradiction certificate.  Second,
all one-generator certificates are factored in closed form.  This proves that information
confined to the $A$-axis or the $B$-axis cannot test compatibility: genuinely mixed minors
are indispensable.  Third, scalar Stieltjes divided differences are lifted to positive
Hankel Gram matrices.  After multiplication by one Vandermonde, these matrices are
positive definite over $\Z$.  Normalizing the underlying Stieltjes measure produces an
\emph{output-sensitive integer lattice}: every nonzero integer polynomial has squared
$L^2$-mass at least the reciprocal of a concrete interpolation determinant.  Fourth, the
same matrices are connected to the Laurent expansion of the rational interpolation
function and hence to Pad\'e and subresultant calculations, giving a fraction-free route
for exact search and formal verification.

The report also derives scale reduction, mixed concavity inequalities, a tail-localization
lemma, integer-Chebyshev contradiction criteria, an exact finite search protocol, and a
Lean-oriented certificate format.  These results do not yet close the conjecture.  The
remaining bottleneck is stated sharply: one must construct mixed semigroup node sets for
which analytic concentration beats the output-dependent lattice threshold, or exhibit a
uniformly wrong mixed subresultant sign for every non-power pair $(A,B)$.
\end{abstract}

\begin{center}
\fbox{\begin{minipage}{0.92\linewidth}\small
\textbf{Status discipline.}
Every assertion below is labelled by type.  ``Proved'' means a complete argument is given.
``Exact experiment'' means integer arithmetic on a finite set, not floating-point evidence.
``Criterion'' is a rigorously sufficient condition whose universal verification remains open.
No claim about an external AI system is used as a mathematical premise, and no unconditional
proof of the \AE\ conjecture is asserted.
\end{minipage}}
\end{center}

\tableofcontents

\section{What is new in this continuation}

The attached report already assembles a broad collection of routes toward the conjecture.
Repeating those reductions would add bulk without adding leverage.  The present continuation
therefore uses only the minimal normal form needed to establish a different interface:
positive integral semigroup data are regarded as samples of the Bernstein function
$u\mapsto u^\theta$, with $0<\theta<1$.  The scalar sign information is then promoted to
quadratic forms, lattices, and subresultants.

The progress ledger is as follows.

\begin{center}
\begin{tabularx}{0.94\linewidth}{@{}p{0.19\linewidth}X p{0.17\linewidth}@{}}
\toprule
Component & Content & Status\\
\midrule
Counterexample normal form & $x>1$ transcendental; $A=2^x$, $B=3^x$, $\theta=1/x\in(0,1)$ & proved\\
Integer sign oracle & Strict signs of mixed generalized Vandermonde determinants & proved\\
Axis factorization & Closed forms on $\{A^j\}$ and $\{B^j\}$; proof that axis-only tests do not couple & proved\\
Moment lattice & Positive definite Hankel matrix becomes integral after one Vandermonde factor & proved\\
Output normalization & Every nonzero $P\in\Z[s]$ has $\|P\|_{L^2(\mu)}^2\ge 1/k_0$ & proved\\
Pad\'e/subresultant bridge & Fraction-free algebraic representation of all moment minors & proved as an algorithmic bridge\\
Tail and Chebyshev tests & Explicit sufficient inequalities for contradiction & conditional criteria\\
Finite search & Exact determinant and moment-minor filters over a bounded candidate box & exact experiment\\
Global conclusion & Uniform mixed certificate for every non-power pair & open\\
\bottomrule
\end{tabularx}
\end{center}

A useful conceptual change is that the arithmetic lower bound is no longer merely
``a nonzero rational has size at least the reciprocal of its denominator.''  The normalized
Stieltjes construction gives a whole positive quadratic form on the coefficient lattice
$\Z^{d+1}$.  The denominator is the integer $k_0$ determined by the \emph{outputs}, and every
nonzero lattice vector has norm at least $1/k_0$.  This is the main new object.

\section{Minimal counterexample normal form}

\begin{proposition}[Elementary reduction]\label{prop:normal}
Assume $2^x,3^x\in\Z$.  Then either $x\in\{0,1\}$ or $x>1$.  If $x\in\Q$, then $x\in\Z$.
Consequently, a nonintegral counterexample must have $x>1$ and must be transcendental.
\end{proposition}

\begin{proof}
If $x<0$, then $0<2^x<1$, so $2^x$ is not an integer.  If $0<x<1$, then
$1<2^x<2$, again impossible.  Now write a positive rational $x=p/q$ in lowest terms.
If $A=2^{p/q}\in\Z$, then $A^q=2^p$.  Taking the $2$-adic valuation gives
$qv_2(A)=p$, hence $q\mid p$ and $q=1$.  Finally, if a nonintegral $x$ were algebraic,
the Gelfond--Schneider theorem would make $2^x$ transcendental, contrary to the hypothesis.
\end{proof}

For the rest of the report, assume hypothetically that a nonintegral counterexample exists
and put
\begin{equation}\label{eq:ABtheta}
 A=2^x\in\N,\qquad B=3^x\in\N,\qquad \theta=\frac1x\in(0,1).
\end{equation}
Then
\begin{equation}\label{eq:reciprocal}
 A^\theta=2,\qquad B^\theta=3.
\end{equation}
The passage to $\theta$ is not cosmetic: all derivatives of $u^\theta$ have a fixed
alternating sign after the first derivative, independently of the size of $x$.

\begin{definition}[Semigroup dictionary]\label{def:dict}
For $e=(m,n)\in\N_0^2$, define
\[
 U_e=A^mB^n,\qquad V_e=2^m3^n.
\]
Then $U_e=V_e^x$ and $V_e=U_e^\theta$.
\end{definition}

Unique factorization makes the $V_e$ distinct.  The same is true of the $U_e$: an equality
$A^mB^n=A^{m'}B^{n'}$ would, after taking the positive $\theta$-th power, give
$2^m3^n=2^{m'}3^{n'}$.  Moreover,
\begin{equation}\label{eq:order}
 U_e<U_{e'}\quad\Longleftrightarrow\quad V_e<V_{e'},
\end{equation}
because $u\mapsto u^\theta$ is strictly increasing.  Thus the order of all input nodes is
known from the corresponding $3$-smooth outputs, without evaluating a logarithm.

\section{A logarithm-free integer sign oracle}

Let $f(u)=u^\theta$.  For distinct positive numbers $u_0<\cdots<u_k$, write
$[u_0,\ldots,u_k]f$ for the ordinary divided difference.

\begin{theorem}[Integral mixed-minor sign hierarchy]\label{thm:signoracle}
Choose distinct exponent vectors $e_0,\ldots,e_k\in\N_0^2$ and order them so that
$V_{e_0}<\cdots<V_{e_k}$.  Set $u_i=U_{e_i}$ and $v_i=V_{e_i}$.  Under the
counterexample hypothesis,
\begin{equation}\label{eq:detcert}
 N(e_0,\ldots,e_k):=
 \det\!\begin{pmatrix}
  1&u_0&u_0^2&\cdots&u_0^{k-1}&v_0\\
  1&u_1&u_1^2&\cdots&u_1^{k-1}&v_1\\
  \vdots&\vdots&\vdots&&\vdots&\vdots\\
  1&u_k&u_k^2&\cdots&u_k^{k-1}&v_k
 \end{pmatrix}
\end{equation}
is a nonzero integer and
\begin{equation}\label{eq:signrule}
 \sgn N(e_0,\ldots,e_k)=(-1)^{k-1}.
\end{equation}
In particular, one wrong sign or one zero determinant is a finite exact certificate that
$(A,B)$ cannot arise from a counterexample.
\end{theorem}

\begin{proof}
The generalized Vandermonde identity gives
\[
 N(e_0,\ldots,e_k)=
 \left(\prod_{0\le i<j\le k}(u_j-u_i)\right)[u_0,\ldots,u_k]f.
\]
The product is positive by \eqref{eq:order}.  The generalized mean-value theorem for divided
differences gives a point $\xi\in(u_0,u_k)$ for which
\[
 [u_0,\ldots,u_k]f=\frac{f^{(k)}(\xi)}{k!}.
\]
For $0<\theta<1$,
\[
 f^{(k)}(u)=\theta(\theta-1)\cdots(\theta-k+1)u^{\theta-k},
\]
whose sign is $(-1)^{k-1}$.  Strictness follows because none of the factors vanishes.
Every entry in \eqref{eq:detcert} is an integer.
\end{proof}

This theorem replaces an equality of logarithmic ratios by an infinite collection of
statements in ordered rings.  It is especially convenient for proof assistants: after the
generic real-analysis lemma is formalized once, each proposed certificate reduces to
normalization of natural-number powers and an integer determinant.

\subsection{Scaling removes a large class of redundant searches}

\begin{proposition}[Semigroup scaling law]\label{prop:scaling}
Let $c=U_e$ and $c^\theta=V_e$.  For any $k+1$ nodes $u_0,\ldots,u_k$ from the semigroup,
\[
 [cu_0,\ldots,cu_k]f=\frac{V_e}{c^k}[u_0,\ldots,u_k]f.
\]
If $N(\UU)$ denotes the determinant numerator in \eqref{eq:detcert}, then
\begin{equation}\label{eq:numscale}
 N(c\UU)=V_e\,c^{k(k-1)/2}N(\UU).
\end{equation}
Hence scaling a certificate by a common semigroup element never changes its sign.
\end{proposition}

\begin{proof}
Homogeneity gives $f(cu)=c^\theta f(u)$, while a divided difference of order $k$ acquires
$c^{-k}$ under dilation of all nodes.  The Vandermonde acquires
$c^{k(k+1)/2}$; multiplying the two scaling factors yields \eqref{eq:numscale}.
\end{proof}

Thus a search may subtract the componentwise minimum of all exponent vectors and retain only
\emph{primitive} patterns touching at least one coordinate axis.  This is a substantial
reduction in exact enumeration and prevents inflated evidence from repeated copies of the
same certificate.

\section{Pure-axis factorization and the necessity of mixed nodes}

The next result is a negative theorem in the useful sense: it identifies a whole family of
methods that cannot couple the two hypotheses.

\begin{theorem}[Closed factorization on one generator]\label{thm:axis}
For every integer $k\ge1$,
\begin{align}
 [1,A,A^2,\ldots,A^k]f
 &=\frac{\displaystyle\prod_{r=0}^{k-1}(2-A^r)}
 {\displaystyle\prod_{r=0}^{k-1}(A^k-A^r)},\label{eq:Aaxis}\\
 [1,B,B^2,\ldots,B^k]f
 &=\frac{\displaystyle\prod_{r=0}^{k-1}(3-B^r)}
 {\displaystyle\prod_{r=0}^{k-1}(B^k-B^r)}.\label{eq:Baxis}
\end{align}
Both expressions automatically have sign $(-1)^{k-1}$ for every $A>2$ and every $B>3$.
Consequently, no argument using only divided differences from one geometric axis can test the
compatibility of $A^\theta=2$ with $B^\theta=3$.
\end{theorem}

\begin{proof}
For the $A$-axis, the determinant numerator has rows indexed by $i=0,\ldots,k$ and columns
$1,A^i,A^{2i},\ldots,A^{(k-1)i},2^i$.  It is therefore the ordinary Vandermonde determinant
in the column bases
\[
 1,A,A^2,\ldots,A^{k-1},2.
\]
Hence it equals
\[
 \left(\prod_{0\le r<s\le k-1}(A^s-A^r)\right)
 \left(\prod_{r=0}^{k-1}(2-A^r)\right).
\]
The node Vandermonde contains the same first product and an additional factor
$\prod_{r=0}^{k-1}(A^k-A^r)$, proving \eqref{eq:Aaxis}.  The $B$-axis is identical with
$2$ replaced by $3$.  The sign count is immediate.
\end{proof}

Formula \eqref{eq:Aaxis} may also be written as a basic-hypergeometric Newton coefficient.
With $q=A^{-1}$, the $k$-th Newton term at an integer $Z$ is
\begin{equation}\label{eq:qnewton}
 A^{-k}\frac{(2;q)_k(Z;q)_k}{(q;q)_k},
 \qquad (a;q)_k=\prod_{r=0}^{k-1}(1-aq^r).
\end{equation}
This exact cancellation explains why a naive denominator estimate on a pure geometric axis
is misleading: most of the generalized Vandermonde has already cancelled.  It also gives a
calibration family against which every proposed mixed estimate should be tested.

\section{The first mixed constraints}

The first nontrivial information appears when the two axes are interlaced.  Since
$1<2<3<4$, the corresponding input nodes satisfy
\[
 1<A<B<A^2.
\]
Strict concavity of $f$ at the first three and the last three nodes gives the following.

\begin{proposition}[Initial rational cone]\label{prop:cone}
Every counterexample pair $(A,B)$ must satisfy
\begin{equation}\label{eq:cone}
 B\ge 2A,\qquad 2B\le A^2+A-1.
\end{equation}
\end{proposition}

\begin{proof}
The consecutive secant slopes at $(1,1),(A,2),(B,3)$ strictly decrease, so
$1/(A-1)>1/(B-A)$ and $B>2A-1$.  Integrality gives $B\ge2A$.  Applying the same argument to
$(A,2),(B,3),(A^2,4)$ gives $1/(B-A)>1/(A^2-B)$, hence $2B<A^2+A$, and integrality gives the
second inequality.
\end{proof}

The order-three sign already adds a nonlinear condition.

\begin{proposition}[First mixed cubic numerator]\label{prop:F3}
Define
\begin{align}
 F_3(A,B)={}&-2A(A-1)^2(A+1)
 +A(A-1)(A+1)(B-1)\notag\\
 &\quad -(A-2)(B-1)(B-A).\label{eq:F3}
\end{align}
Then every counterexample pair satisfies $F_3(A,B)>0$.
\end{proposition}

\begin{proof}
Let $P_2$ be the quadratic interpolant to $f$ at $1,A,A^2$.  Newton interpolation and
\eqref{eq:Aaxis} give
\[
 P_2(z)=1+\frac{z-1}{A-1}
 +\frac{(2-A)(z-1)(z-A)}{A(A-1)^2(A+1)}.
\]
The interpolation error at $B\in(A,A^2)$ is
\[
 3-P_2(B)=[1,A,B,A^2]f\,(B-1)(B-A)(B-A^2).
\]
The divided difference is positive, while the product is negative.  Thus $P_2(B)>3$.
Clearing the positive denominator $A(A-1)^2(A+1)$ gives \eqref{eq:F3}.
\end{proof}

Higher canonical constraints are obtained from the ordered $3$-smooth sequence
\[
 1,2,3,4,6,8,9,12,16,18,24,27,32,36,\ldots
\]
by replacing each $2^m3^n$ with $A^mB^n$ and imposing \eqref{eq:signrule} on every selected
subset.  The result is a nested sequence of semialgebraic regions in the integer lattice.
Unlike a numerical comparison of logarithms, every exclusion comes with a small list of
exponent pairs and an integer determinant that can be checked independently.

\section{From scalar Stieltjes differences to positive Gram matrices}

The function $f(u)=u^\theta$, $0<\theta<1$, is a complete Bernstein function.  Its elementary
Stieltjes representation is
\begin{equation}\label{eq:stieltjes}
 u^\theta=\cst\int_0^\infty \frac{u}{u+s}s^{\theta-1}\dd s,
 \qquad \cst=\frac{\sin(\pi\theta)}{\pi}.
\end{equation}
For $n\ge1$ and distinct positive nodes $u_0,\ldots,u_n$, taking an order-$n$ divided
difference under the integral gives
\begin{equation}\label{eq:stieltjesdd}
 (-1)^{n-1}[u_0,\ldots,u_n]f
 =\cst\int_0^\infty\frac{s^\theta}{\prod_{i=0}^n(s+u_i)}\dd s>0.
\end{equation}
The attached report uses scalar divided-difference positivity.  The new step here is to
insert polynomial squares into the same measure and keep exact control of their arithmetic.

Fix an ordered semigroup node set
\[
 \UU=(u_0<\cdots<u_n),\qquad v_i=u_i^\theta\in\N,
\]
and an integer $d\ge0$ with
\begin{equation}\label{eq:degreecondition}
 n\ge2d+1.
\end{equation}
For $0\le r\le2d$, define
\begin{equation}\label{eq:hr}
 h_r(\UU)=\cst\int_0^\infty
 \frac{s^{\theta+r}}{\prod_{i=0}^n(s+u_i)}\dd s.
\end{equation}
Condition \eqref{eq:degreecondition} guarantees convergence at infinity.

\begin{lemma}[Rational moment formula]\label{lem:rationalmoment}
Let $P_\UU(z)=\prod_{i=0}^n(z-u_i)$.  Then
\begin{equation}\label{eq:momentformula}
 h_r(\UU)=(-1)^{n+r+1}
 \sum_{i=0}^n\frac{u_i^rv_i}{P_\UU'(u_i)}
 =(-1)^{n+r+1}[u_0,\ldots,u_n]z^{\theta+r}.
\end{equation}
In particular, every $h_r(\UU)$ is rational and positive.
\end{lemma}

\begin{proof}
For a complex parameter $\beta$ in the convergence strip $-1<\Re\beta<n$, consider
\[
 I(\beta)=\int_0^\infty\frac{s^\beta}{\prod_i(s+u_i)}\dd s.
\]
In the substrip $-1<\Re\beta<0$, partial fractions may be integrated term by term, using
\[
 \int_0^\infty\frac{s^\beta}{s+u}\dd s=-\frac{\pi u^\beta}{\sin\pi\beta}.
\]
Since the coefficient of $(s+u_i)^{-1}$ is $(-1)^n/P_\UU'(u_i)$, one obtains
\[
 I(\beta)=(-1)^{n+1}\frac{\pi}{\sin\pi\beta}
 \sum_i\frac{u_i^\beta}{P_\UU'(u_i)}.
\]
Both sides continue analytically through the full convergence strip, with removable
singularities at integral parameters after the sum is formed.  Put $\beta=\theta+r$ and use
$\sin\pi(\theta+r)=(-1)^r\sin\pi\theta$.
\end{proof}

\section{The output-dependent moment lattice}

Let
\begin{equation}\label{eq:Delta}
 \Delta(\UU)=\prod_{0\le i<j\le n}(u_j-u_i)>0,
 \qquad k_r(\UU)=\Delta(\UU)h_r(\UU).
\end{equation}

\begin{theorem}[Integral positive moment matrix]\label{thm:momentlattice}
Under the counterexample hypothesis and \eqref{eq:degreecondition}, every $k_r(\UU)$ is a
positive integer.  The Hankel matrix
\begin{equation}\label{eq:Kmatrix}
 K_d(\UU)=\bigl(k_{i+j}(\UU)\bigr)_{0\le i,j\le d}
\end{equation}
is an integral positive-definite matrix.  Consequently,
\begin{equation}\label{eq:detlower}
 \det K_d(\UU)\in\N,
 \qquad
 \det H_d(\UU)\ge\Delta(\UU)^{-(d+1)},
\end{equation}
where $H_d=(h_{i+j})$.
\end{theorem}

\begin{proof}
By \eqref{eq:momentformula},
\[
 k_r=(-1)^{n+r+1}\sum_i u_i^rv_i\frac{\Delta(\UU)}{P_\UU'(u_i)}.
\]
The quotient $\Delta(\UU)/P_\UU'(u_i)$ is an integer: after cancellation it is, up to sign,
the product of all pairwise differences not incident with $u_i$.  Positivity follows from
the integral definition \eqref{eq:hr}.

For a nonzero vector $c=(c_0,\ldots,c_d)^T\in\R^{d+1}$ and
$Q_c(s)=\sum_{j=0}^dc_js^j$,
\[
 c^TK_dc=\Delta(\UU)\cst\int_0^\infty
 \frac{s^\theta Q_c(s)^2}{\prod_i(s+u_i)}\dd s>0.
\]
Thus $K_d$ is positive definite.  Its determinant is a positive integer, so it is at least
one.  Since $K_d=\Delta(\UU)H_d$, \eqref{eq:detlower} follows.
\end{proof}

The normalization reveals the promised output sensitivity.  Put
\begin{equation}\label{eq:k0}
 k_0(\UU)=\Delta(\UU)h_0(\UU)
 =\left|\det\begin{pmatrix}
 1&u_0&\cdots&u_0^{n-1}&v_0\\
 \vdots&\vdots&&\vdots&\vdots\\
 1&u_n&\cdots&u_n^{n-1}&v_n
 \end{pmatrix}\right|\in\N
\end{equation}
and define the probability measure
\begin{equation}\label{eq:mu}
 \dd\mu_\UU(s)=\frac{\cst}{h_0(\UU)}
 \frac{s^\theta}{\prod_{i=0}^n(s+u_i)}\dd s.
\end{equation}

\begin{corollary}[Output-dependent $L^2$ lattice]\label{cor:L2lattice}
For every nonzero polynomial $Q\in\Z[s]$ of degree at most $d$,
\begin{equation}\label{eq:L2lower}
 \int_0^\infty Q(s)^2\dd\mu_\UU(s)
 =\frac{M_\UU(Q)}{k_0(\UU)}\ge\frac1{k_0(\UU)},
\end{equation}
where $M_\UU(Q)$ is a positive integer.  In addition, if
$G_d=(h_{i+j}/h_0)_{0\le i,j\le d}$ is the normalized Gram matrix, then
\begin{equation}\label{eq:Gdet}
 \det G_d(\UU)=\frac{\det K_d(\UU)}{k_0(\UU)^{d+1}}
 \ge k_0(\UU)^{-(d+1)}.
\end{equation}
\end{corollary}

\begin{proof}
If $c\in\Z^{d+1}\setminus\{0\}$ is the coefficient vector of $Q$, then
$c^TK_dc$ is a positive integer.  Divide by $k_0=\Delta h_0$.  The determinant identity is
immediate.
\end{proof}

The denominator in \eqref{eq:L2lower} is not a generic common denominator.  It is the
interpolation determinant formed from the actual output values $2^m3^n$.  This is the point
at which cancellation specific to the conjecture enters the moment method.

\subsection{Unimodular recentering}

For an integer $C$, the bases $1,s,\ldots,s^d$ and
$1,(s-C),\ldots,(s-C)^d$ are related by a unitriangular integral matrix.  Therefore
\begin{equation}\label{eq:centered}
 \det\bigl(\mathbb E_\mu[(S-C)^{i+j}]\bigr)_{0\le i,j\le d}
 =\det G_d(\UU).
\end{equation}
This permits analytic estimates around an integer center without weakening the arithmetic
lower bound.  More generally, every unimodular change of polynomial basis preserves the
lattice determinant.

\subsection{Smith and divisibility refinements}

The lower bound ``at least one'' for $\det K_d$ is deliberately conservative.  In exact
examples, the matrix often has large systematic divisors.  A fraction-free Smith-normal-form
calculation yields invariant factors
\[
 s_1\mid s_2\mid\cdots\mid s_{d+1},
 \qquad \det K_d=s_1\cdots s_{d+1}.
\]
Any uniform divisor extracted from the mixed output pattern strengthens \eqref{eq:Gdet}.
This is one of the most concrete current targets: pure-axis divisors factor completely by
\cref{thm:axis}, while mixed invariant factors can detect compatibility.

\section{Interpolation rational functions and subresultants}

Let $R_\UU(z)$ be the unique polynomial of degree at most $n$ satisfying
$R_\UU(u_i)=v_i$.  The partial fraction expansion gives
\begin{equation}\label{eq:RP}
 \frac{R_\UU(z)}{P_\UU(z)}
 =\sum_{i=0}^n\frac{v_i}{P_\UU'(u_i)(z-u_i)}
 =\sum_{r\ge0}m_rz^{-r-1},
\end{equation}
where
\begin{equation}\label{eq:mr}
 m_r=\sum_{i=0}^n\frac{u_i^rv_i}{P_\UU'(u_i)},
 \qquad h_r=(-1)^{n+r+1}m_r.
\end{equation}
Thus the positive moment matrices are signed Hankel blocks of the Laurent expansion at
infinity of a rational interpolation function determined entirely by integer input-output
data.

\begin{proposition}[Fraction-free Pad\'e bridge]\label{prop:pade}
After clearing the single factor $\Delta(\UU)$, every Hankel minor needed in
\cref{thm:momentlattice} can be computed in either of two equivalent exact ways:
\begin{enumerate}[label=(\roman*)]
\item from the integral-moment formula \eqref{eq:momentformula} followed by a Bareiss
      determinant;
\item from the Pad\'e table at infinity of $R_\UU/P_\UU$, equivalently from the signed
      subresultant polynomial-remainder sequence of $P_\UU$ and $R_\UU$.
\end{enumerate}
Nonvanishing of the relevant Hankel determinants is equivalent to normality of the
corresponding Pad\'e entries.  Under a counterexample, the convention-adjusted signs are
fixed by positivity.
\end{proposition}

\begin{proof}[Proof sketch]
Equation \eqref{eq:RP} identifies the Laurent coefficients.  The denominator of a Pad\'e
approximant is obtained by solving a Hankel linear system in these coefficients; its
solvability and degree are controlled by consecutive Hankel minors.  The classical
Hankel--subresultant identity identifies the same systems with principal minors in the
Sylvester/subresultant construction for $P_\UU$ and $R_\UU$.  Because $P_\UU$ is monic and
$\Delta(\UU)R_\UU\in\Z[z]$, fraction-free elimination remains in $\Z$ throughout.  See
\cite{wall,bakergraves,krattenthaler} for the standard identities.
\end{proof}

This bridge matters for two reasons.  First, subresultant signs can be certified without
approximating $\theta$.  Second, a proof assistant need not formalize all of Pad\'e theory to
check a discovered certificate: the final object is simply an integer determinant or a
fraction-free remainder with a specified sign.

\section{Concrete contradiction criteria}

The moment lattice converts the conjecture into a competition between an analytic upper
bound and an arithmetic lower bound.

\begin{criterion}[Single-polynomial certificate]\label{crit:poly}
A counterexample is impossible if one can find a semigroup node set $\UU$ with
$n\ge2d+1$ and a nonzero $Q\in\Z[s]$, $\deg Q\le d$, such that
\begin{equation}\label{eq:polycrit}
 \cst\int_0^\infty
 \frac{s^\theta Q(s)^2}{\prod_i(s+u_i)}\dd s
 <\frac1{\Delta(\UU)}.
\end{equation}
Equivalently, it is enough to prove
$\int Q^2\dd\mu_\UU<1/k_0(\UU)$.
\end{criterion}

This is immediate from \cref{thm:momentlattice,cor:L2lattice}, but it is useful because the
left side can be bounded without expanding the enormous determinant $k_0$.

\begin{criterion}[Determinant certificate]\label{crit:det}
A counterexample is impossible if, for some $\UU$ and $d$,
\begin{equation}\label{eq:detcrit}
 \det\bigl(\mathbb E_\mu[(S-C)^{i+j}]\bigr)_{0\le i,j\le d}
 <k_0(\UU)^{-(d+1)}
\end{equation}
for one (hence every) integer center $C$.
By Hadamard's inequality it is enough that
\begin{equation}\label{eq:hadamardcrit}
 \prod_{j=0}^d\mathbb E_\mu[(S-C)^{2j}]<k_0(\UU)^{-(d+1)}.
\end{equation}
\end{criterion}

The determinant version can succeed even when no explicit short polynomial is guessed: a
geometry-of-numbers or LLL step can then recover one.  Conversely, the exact integral matrix
$K_d$ provides a direct check that a numerical short vector is not an artefact of rounding.

\subsection{A tail-localization lemma}

Write
\[
 W_\UU(s)=\frac{s^\theta}{\prod_i(s+u_i)},
 \qquad N_\UU(s)=\#\{i:u_i\le s\}.
\]

\begin{lemma}[Node-count tail contraction]\label{lem:tail}
For $\lambda\ge1$ and $s>0$,
\begin{equation}\label{eq:tailpoint}
 \frac{W_\UU(\lambda s)}{W_\UU(s)}
 \le \lambda^\theta
 \left(\frac{2}{\lambda+1}\right)^{N_\UU(s)}.
\end{equation}
Consequently, for every $R>0$,
\begin{equation}\label{eq:tailinterval}
 \int_{\lambda R}^{\lambda^2R}W_\UU(s)\dd s
 \le \lambda^{\theta+1}
 \left(\frac{2}{\lambda+1}\right)^{N_\UU(R)}
 \int_R^{\lambda R}W_\UU(s)\dd s.
\end{equation}
\end{lemma}

\begin{proof}
The ratio is
\[
 \lambda^\theta\prod_i\frac{s+u_i}{\lambda s+u_i}.
\]
If $u_i\le s$, put $t=u_i/s\in(0,1]$.  Then
$(1+t)/(\lambda+t)\le2/(\lambda+1)$; the remaining factors are at most one.
For \eqref{eq:tailinterval}, substitute $s=\lambda t$ and use
$N_\UU(t)\ge N_\UU(R)$ on $[R,\lambda R]$.
\end{proof}

For a node set consisting of an initial segment of the mixed semigroup, $N_\UU(s)$ grows
stepwise with the number of available $3$-smooth values.  Iterating \eqref{eq:tailinterval}
therefore gives rigorous tail bounds that improve as additional mixed nodes are inserted.
The unresolved comparison is whether this concentration can outrun the simultaneous growth
of $k_0(\UU)$.

\subsection{An integer-Chebyshev target on a unit interval}

For an integer $C$, the polynomial
\[
 T_C(s)=(s-C)(s-C-1)
\]
has integer coefficients and $|T_C(s)|\le1/4$ on $[C,C+1]$.  Hence
$Q_m=T_C^m$ obeys $|Q_m|^2\le16^{-m}$ there.  Splitting the normalized integral gives
\begin{equation}\label{eq:unitinterval}
 \|Q_m\|_{L^2(\mu_\UU)}^2
 \le16^{-m}+
 \int_{[0,\infty)\setminus[C,C+1]}|T_C(s)|^{2m}\dd\mu_\UU(s).
\end{equation}
Combining \eqref{eq:unitinterval} with \cref{lem:tail} gives a fully explicit route to
\cref{crit:poly}.  It is not enough to show mere concentration: the tail must be controlled
with the polynomial growth included, and the resulting bound must beat $1/k_0$.

\begin{problem}[Quantitative localization problem]\label{prob:localize}
Construct, from every hypothetical non-power pair $(A,B)$, primitive mixed node sets
$\UU_j$, integer intervals $[C_j,C_j+1]$, and degrees $m_j$ for which the right side of
\eqref{eq:unitinterval} is $o(1/k_0(\UU_j))$.
\end{problem}

Integral-exponent controls $(A,B)=(2^r,3^r)$ show that a bound depending only on crude node
sizes cannot solve \cref{prob:localize}; it must use mixed arithmetic, for example Smith
invariant factors or unusually small mixed subresultants.

\section{Exact finite experiments}

All computations in this section use Python integers, the fraction-free Bareiss determinant,
and exact exponentiation.  No logarithm is used in a certificate.  Floating point is used
only to select a few illustrative near-power pairs for display; their acceptance or rejection
is then decided exactly.

\subsection{Bounded mixed-sign hierarchy}

The experiment enumerates
\[
 3\le A\le60,\qquad 2A\le B\le\frac{A^2+A-1}{2},
\]
then imposes order agreement on the first twelve $3$-smooth outputs and all subset determinant
signs through a fixed order.  Common-power controls are pairs $(2^r,3^r)$ in the box.

\begin{center}
\emph{The exact enumeration fragment was unavailable at build time.}
\end{center}

\noindent\textit{Machine-generated summary.} The computation is an exact finite experiment; no global conclusion is inferred from it.

The experiment is deliberately interpreted conservatively.  A surviving non-control pair is
not evidence for a solution; it only means that the tested finite sign fragment has not yet
located a contradiction.  An excluded pair, by contrast, comes with a finite exact witness.

\subsection{Sample exact certificates}

For each displayed pair, the program first checks semigroup order and then searches all
subsets of the first fourteen smooth outputs through order six.

\begin{center}
\emph{The sample-certificate fragment was unavailable at build time.}
\end{center}

The pure-power controls are an important guard against sign-convention errors.  They must pass
every test because they correspond to the genuine exponent $x\in\Z$.

\subsection{Moment-minor filters}

Using the first eight smooth outputs
$1,2,3,4,6,8,9,12$, one has $n=7$ and may build the $4\times4$ matrix $K_3$.
The following table reports the signs of its leading principal minors.

\begin{center}
\emph{The moment-matrix fragment was unavailable at build time.}
\end{center}

A nonpositive leading minor is a compact contradiction certificate.  Positive leading minors
through this size do not establish the existence of a power exponent; they merely show that
this particular truncated Stieltjes moment test has not failed.

\subsection{Reproducible pseudocode}

\begin{verbatim}
for each integer pair (A,B) in the rational cone:
    E := exponent pairs of the first N numbers 2^m 3^n
    U[e] := A^m B^n;  V[e] := 2^m 3^n
    reject if U is not ordered as V
    for k = 2,...,K:
        for every selected (k+1)-subset I of E:
            D := det( 1, U, ..., U^(k-1), V ) on I
            reject unless sign(D) = (-1)^(k-1)

for a moment certificate with n+1 nodes and n >= 2d+1:
    Delta := product_{i<j} (U[j]-U[i])
    k[r] := (-1)^(n+r+1) sum_i U[i]^r V[i] Delta/P'(U[i])
    K := Hankel(k[0],...,k[2d])
    reject unless all required principal minors of K are positive
\end{verbatim}

The determinant numerator itself, together with the exponent vectors, is a portable
certificate.  A checker does not need the search code.

\section{A third-base bridge and why it remains attractive}

The Six Exponentials Theorem gives a clean strategic target.  If $x$ is irrational, then
$1,x$ are linearly independent over $\Q$, while $\log2,\log3,\log5$ are linearly independent
over $\Q$.  Therefore at least one of
\[
 2,3,5,2^x,3^x,5^x
\]
is transcendental.  Under a counterexample, the first five numbers except $5^x$ are algebraic,
so the theorem forces $5^x$ to be transcendental.  Consequently:

\begin{criterion}[Third-base bootstrap]\label{crit:thirdbase}
Any argument deriving the algebraicity of $5^x$ from $2^x,3^x\in\Z$ proves the \AE\
conjecture.  The same holds with $5$ replaced by any algebraic number whose logarithm is
$\Q$-independent of $\log2$ and $\log3$.
\end{criterion}

The moment framework suggests one possible bridge: interpolate $y\mapsto y^x$ at many
$3$-smooth integers, where all values are integral, and evaluate at $5$.  What is missing is
an arithmetic approximation strong enough to force algebraicity, not merely a good real
approximation.  Ordinary Lagrange denominators grow too rapidly.  A successful construction
would likely need one of the output-specific cancellations isolated in \cref{thm:axis,
cor:L2lattice,prop:pade}.

\section{Routes tested in this continuation and the precise obstruction}

This section records three tempting diversions so that subsequent work does not restart them
without addressing the stated obstruction.

\subsection{Dense rational graph}

If a counterexample exists, then for every $m,n\in\Z$,
\[
 (2^{m+nx},3^{m+nx})=(2^mA^n,3^mB^n)\in\Q_{>0}^2.
\]
Thus the real analytic curve $t\mapsto(2^t,3^t)$ contains a rank-two subgroup of rational
points that is dense along the real one-parameter subgroup.  This is an appealing
functional-transcendence picture, but a direct Pila--Wilkie count is too weak: points in a
fixed compact parameter interval obtained with $|m|,|n|\le N$ have height exponential in
$N$ and number only $O(N)$, i.e. logarithmic in height.  This lies far below a generic
$H^\varepsilon$ upper bound.  A useful return to this route would require a sharp
polylogarithmic count specific to irrational power curves.

\subsection{$p$-adic automatic continuity}

The group homomorphism generated by $2\mapsto3$ and $A\mapsto B$ is real-order preserving.
It is tempting to extend it continuously to a $p$-adic closure and compare $p$-adic logarithm
ratios.  The missing implication is continuity: a congruence
$2^mA^n\equiv1\pmod{p^r}$ gives no known congruence for $3^mB^n$.  Real closeness and
$p$-adic closeness pull in unrelated directions.  A future use of this idea needs an
independent height-to-congruence transfer principle.

\subsection{Arithmetic progressions of smooth bases}

The values of $y^x$ are integral at every $3$-smooth $y$.  Primitive arithmetic progressions
such as $1,2,3,4$ and $2,9,16$ yield exact convexity inequalities in $A$ and $B$, but known
$S$-unit finiteness prevents an arbitrarily long supply of such stencils.  Their best role is
as low-cost filters and seeds for mixed determinants, not as a stand-alone route to a third
base.

\section{Formalization blueprint}

The sign certificates were designed to separate analysis from computation.

\begin{enumerate}[label=\textbf{Layer \arabic*.},wide]
\item \textbf{Generic analysis.}  Formalize that for $0<\theta<1$, the order-$k$ divided
      difference of $u^\theta$ on increasing positive nodes has sign $(-1)^{k-1}$.
      A generalized mean-value theorem is sufficient; no Stieltjes integration is needed for
      the first certificate checker.
\item \textbf{Semigroup specialization.}  From $A^\theta=2$ and $B^\theta=3$, prove
      $(A^mB^n)^\theta=2^m3^n$ and transfer output order to input order.
\item \textbf{Integer determinant.}  Identify the divided-difference numerator with the
      generalized Vandermonde determinant and cast its entries into $\Z$.
\item \textbf{Certificate evaluation.}  For a fixed exponent list, use kernel reflection or
      a fraction-free determinant tactic to prove the explicit sign.
\item \textbf{Moment extension.}  Separately formalize the beta-integral identity in
      \cref{lem:rationalmoment}, then derive integral positive definiteness of $K_d$.
\end{enumerate}

A schematic Lean data structure is:
\begin{verbatim}
structure ExpNode where
  m : Nat
  n : Nat

def inputNode  (A B : Nat) (e : ExpNode) := A^e.m * B^e.n
def outputNode (e : ExpNode) := 2^e.m * 3^e.n

def certificateNumerator (A B : Nat) (E : Fin (k+1) -> ExpNode) : Int :=
  Matrix.det (fun i j =>
    if j < k then (inputNode A B (E i) : Int)^j
    else (outputNode (E i) : Int))
\end{verbatim}
The final contradiction theorem takes as input the sortedness proof and a normalized integer
value for the determinant.  This makes discovered certificates auditable even before the
full search is formalized.

\section{Progress assessment and next attacks}

The new framework changes the main question from ``How can one estimate another scalar
divided difference?'' to two more structured problems.

\begin{problem}[Uniform mixed-sign problem]\label{prob:mixedsign}
Show that for every integer pair $(A,B)$ satisfying the elementary order constraints but not
of the form $(2^r,3^r)$, some primitive finite exponent pattern has a determinant sign
contrary to \eqref{eq:signrule}; preferably with an order bounded effectively in terms of
$\log A+\log B$.
\end{problem}

\begin{problem}[Moment-lattice gap problem]\label{prob:gap}
Construct mixed node sets for which either a principal minor of $K_d$ is nonpositive, or an
analytic estimate proves $\det K_d<1$.  Equivalently, find an integer polynomial violating
\eqref{eq:L2lower} by a provable Stieltjes concentration estimate.
\end{problem}

The immediate research program is:
\begin{enumerate}[label=(\arabic*)]
\item enumerate primitive exponent patterns modulo the scaling law \eqref{eq:numscale};
\item factor low-order mixed determinant numerators symbolically and identify divisibility
      that distinguishes common-power controls;
\item compute Smith invariants and subresultant remainder signs for initial smooth-node sets;
\item optimize integer-Chebyshev polynomials against the exact normalized measure, with
      rigorous tail bounds from \cref{lem:tail};
\item convert any recurring negative minor into a small formal certificate and search for a
      parameter-uniform factorization.
\end{enumerate}

The strongest proved advance is \cref{thm:momentlattice,cor:L2lattice}: a hypothetical
counterexample must support, for every admissible mixed node set, a positive definite
\emph{integral} Hankel form whose normalized coefficient lattice has a quantified minimum.
The strongest negative advance is \cref{thm:axis}: no amount of refinement confined to a
single generator can settle compatibility.  Together they sharply focus subsequent effort
on output-sensitive mixed minors and their analytic concentration.

\appendix

\section{A compact dependency graph of the main results}

\begin{center}
\begin{tabularx}{0.94\linewidth}{@{}p{0.26\linewidth}X@{}}
\toprule
Result & Dependencies\\
\midrule
\cref{thm:signoracle} & reciprocal normal form; generalized mean-value theorem; generalized Vandermonde identity\\
\cref{prop:scaling} & homogeneity of $u^\theta$ and Vandermonde scaling\\
\cref{thm:axis} & standard Vandermonde determinant in the column bases $1,A,\ldots,A^{k-1},2$\\
\cref{lem:rationalmoment} & beta integral; partial fractions; analytic continuation within the convergence strip\\
\cref{thm:momentlattice} & \cref{lem:rationalmoment}; divisibility of $\Delta/P'(u_i)$; positivity of a square integral\\
\cref{cor:L2lattice} & integral positive definiteness and normalization by $k_0$\\
\cref{prop:pade} & Laurent expansion of $R/P$; classical Pad\'e--Hankel--subresultant identities\\
\cref{lem:tail} & pointwise comparison of the Stieltjes weight under dilation\\
\bottomrule
\end{tabularx}
\end{center}

\section{Sanity checks against integral exponents}

For $x=r\in\N$, one has $A=2^r$, $B=3^r$, and $\theta=1/r$.  Every theorem above must remain
consistent with these controls.  In particular:
\begin{itemize}
\item the determinant signs in \cref{thm:signoracle} hold strictly;
\item the axis factorizations are exact;
\item every $K_d$ is integral positive definite;
\item no valid analytic estimate can make $\det K_d<1$;
\item any proposed ``universal'' localization bound that also excludes these controls has
      lost essential mixed arithmetic information.
\end{itemize}
This control family is built into the exact search and is a practical defense against an
incorrect sign, a missing Vandermonde factor, or a non-rigorous numerical Cholesky step.

\section{Bibliographic orientation}

The conjecture originates in the work of Alaoglu and Erd\H{o}s on highly composite and
related numbers \cite{alaogluerdos}.  Its transcendence-theoretic setting is the Four
Exponentials Conjecture and the proved Six Exponentials Theorem; standard references include
\cite{lang,waldschmidt,baker}.  The Gelfond--Schneider reduction used in
\cref{prop:normal} is classical \cite{gelfond,baker}.  Bernstein and Stieltjes representations
are treated systematically in \cite{schilling,akhiezer}; total positivity and sign-regular
kernels in \cite{karlin}.  The $q$-Newton notation in \eqref{eq:qnewton} follows
\cite{gasperrahman}.  Pad\'e, continued-fraction, determinant, and subresultant techniques are
covered in \cite{wall,bakergraves,krattenthaler}.  Exact real-algebraic and fraction-free
certificate methodology is consonant with \cite{basupollackroy}.

\begin{thebibliography}{99}
\bibitem{alaogluerdos}
L. Alaoglu and P. Erd\H{o}s,
\emph{On highly composite and similar numbers},
Transactions of the American Mathematical Society \textbf{56} (1944), 448--469.

\bibitem{akhiezer}
N. I. Akhiezer,
\emph{The Classical Moment Problem and Some Related Questions in Analysis},
Oliver \& Boyd, 1965.

\bibitem{baker}
A. Baker,
\emph{Transcendental Number Theory}, second edition,
Cambridge University Press, 1990.

\bibitem{bakergraves}
G. A. Baker, Jr. and P. Graves-Morris,
\emph{Pad\'e Approximants}, second edition,
Cambridge University Press, 1996.

\bibitem{basupollackroy}
S. Basu, R. Pollack, and M.-F. Roy,
\emph{Algorithms in Real Algebraic Geometry}, second edition,
Springer, 2006.

\bibitem{gasperrahman}
G. Gasper and M. Rahman,
\emph{Basic Hypergeometric Series}, second edition,
Cambridge University Press, 2004.

\bibitem{gelfond}
A. O. Gelfond,
\emph{Transcendental and Algebraic Numbers},
Dover, 1960.

\bibitem{karlin}
S. Karlin,
\emph{Total Positivity}, volume I,
Stanford University Press, 1968.

\bibitem{krattenthaler}
C. Krattenthaler,
\emph{Advanced determinant calculus},
S\'eminaire Lotharingien de Combinatoire \textbf{42} (1999), Article B42q.

\bibitem{lang}
S. Lang,
\emph{Introduction to Transcendental Numbers},
Addison--Wesley, 1966.

\bibitem{schilling}
R. L. Schilling, R. Song, and Z. Vondra\v{c}ek,
\emph{Bernstein Functions: Theory and Applications}, second edition,
De Gruyter, 2012.

\bibitem{waldschmidt}
M. Waldschmidt,
\emph{Diophantine Approximation on Linear Algebraic Groups},
Springer, 2000.

\bibitem{wall}
H. S. Wall,
\emph{Analytic Theory of Continued Fractions},
Van Nostrand, 1948; reprint, AMS Chelsea.
\end{thebibliography}

\end{document}
